\section{Lesson 2: Multi-Member Districts Partially but Not Fully Escape Geographic Sorting}
\label{sec:lesson2}

If single-member geographic districts cannot escape the Rodden sorting gap,
can multi-member districts?  Paper E.1 provides a detailed answer: partially,
but not fully, and the degree of escape increases with district size in a
predictable way \citep{dellaLibera2026e1}.

\subsection{The Mechanism of Partial Escape}

The mechanism is straightforward.  In a single-member district centered on
an urban core, a 70--30 Democratic advantage produces one Democratic seat and
zero Republican seats.  In a three-member district centered on the same core
with a surrounding suburban ring, the overall composition might shift to 60--40
Democratic, producing two Democratic seats and one Republican seat under
proportional allocation.  The margin improvement from 70--30 to 60--40 is
modest, but the seat change from 1--0 to 2--1 is a 33\% gain in Republican
representation --- a meaningful partial escape from sorting-driven
underrepresentation.

The key variable is how much suburban and exurban population can be incorporated
into an enlarged district without exceeding the population quota.  Three-member
districts have three times the population of single-member districts, so they
reach further into suburban rings; five-member districts reach further still.
The extent to which this averaging across geographic zones reduces partisan
concentration depends on the city-suburb-rural gradient of the specific
metropolitan area.

\subsection{Empirical Evidence from E.1}

Paper E.1 implements three-member and five-member districts using the same
edge-weighted recursive bisection algorithm as the baseline but with population
quotas scaled by district size \citep{dellaLibera2026e1}.  The 50-state results
show:

\begin{itemize}
  \item \textbf{Three-member districts}: Seat--vote deviation reduced from
        the median sorting gap of $+4.2$~pp Republican (range: 1.1~pp in ND
        to 11.3~pp in WI; baseline) to $+2.5$~pp Republican, a 40\%
        reduction in the sorting-driven gap.  The gap reduction is smallest
        in low-sorting states (IA, ND) and largest in high-sorting states
        (IL, NY, PA, MI).  Mean Polsby--Popper compactness declines from
        0.367 to 0.286, a 22\% reduction.

  \item \textbf{Five-member districts}: Seat--vote deviation reduced to
        $+1.5$~pp Republican, a 64\% reduction.  Mean Polsby--Popper
        compactness declines to 0.241, a 34\% reduction.
\end{itemize}

The compactness decline is expected: larger districts must span more
geographic territory and more heterogeneous terrain, producing less regular
shapes.  The important question is whether the proportionality improvement
is worth the compactness cost.

\subsection{Why Full Escape Is Impossible}

The residual $+1.5$~pp gap that persists even with five-member districts
reflects extreme urban concentration in a small number of metropolitan areas.
Manhattan (New York County) votes approximately 85--15 Democratic in
congressional elections; central Chicago similarly at 80--20; central Los
Angeles at 75--25.  These concentrations are so extreme that no realistically
sized district can incorporate enough surrounding population to reduce the
Democratic margin to proportional territory.

A five-member district centered on Manhattan with a population quota five times
the single-member benchmark would encompass much of the outer boroughs, but
the outer boroughs themselves are heavily Democratic, so the district-level
composition remains far above proportional levels.  The algorithm cannot
extend the district into New Jersey or Connecticut without violating the state-
boundary constraint.

The theoretical maximum reduction in the sorting gap from multi-member
districting, holding state boundaries constant and allowing unlimited district
size, is approximately 80\% of the current gap.  The remaining 20\% is
irreducible given the intensity of concentration in the largest American
cities.\footnote{The 20\% irreducibility bound is estimated as follows: in the
most geographically sorted state (WI, sorting gap $= 11.3$~pp), the minimum
achievable partisan gap under any contiguous equal-population partition of 2020
census tracts is approximately 2.3~pp \citep{herschlag2020quantifying} (NC
ensemble lower bound, scaled by the WI/NC sorting ratio).  This represents
20\% of the 11.3~pp total gap.  The remaining 80\% is structurally addressable
through algorithm design, as demonstrated by the E.1 multi-member results.}
This bound is derived from the fraction of Democratic votes that are in census
tracts with $>70\%$ Democratic composition --- the ``captive'' Democratic votes
that no district boundary can dilute.

\subsection{The Size--Proportionality Trade-off}

The data suggest a concave relationship across the three tested configurations
(single-member, 3-member, 5-member districts), consistent with diminishing
returns to increasing district magnitude: the marginal gap reduction is
approximately 40\% for the first additional two members (one to three) and
approximately 24\% for the next two (three to five).  With only three data
points, the functional form cannot be confirmed, but the pattern is sufficient
to support the policy conclusion that the Fair Representation Act's 3--5 member
specification captures most of the available improvement.  The marginal gain
diminishes because the most easily diversifiable geographic zones are
incorporated first --- the suburban rings adjacent to urban cores --- and what
remains is the irreducible urban concentration.

This concavity has practical implications.  The Fair Representation Act's
choice of three-to-five-member districts is near the knee of the
marginal-benefit curve: it captures most of the available proportionality
improvement while keeping ballot complexity and district-size effects manageable.
Larger districts (seven or more members) buy little additional proportionality
at substantially higher institutional cost.

\subsection{Minority Representation Benefits}

Multi-member districts provide a second proportionality benefit that
single-member districts cannot: minority representation that is
proportional to minority population share rather than dependent on creating
majority-minority districts.  In a single-member system, the \emph{Gingles}
framework \citep{gingles1986} requires majority-minority districts as a VRA
remedy, but these districts can only be created where minority populations
are geographically concentrated enough to form a majority.

In a three-member district, a 35\% minority population share translates to
one minority seat if minority voters prefer the minority-group candidate and
bloc-vote accordingly --- without requiring a majority-minority geographic
concentration.  Paper E.1 finds that three-member districts produce
approximately 18\% more effective minority representation than the baseline
single-member system when minority populations are moderately dispersed
\citep{dellaLibera2026e1}.  This benefit complements rather than duplicates
the VRA majority-minority district guarantee.
